% Requires BibTeX entries: cnil2024reuse, strubell2019energy,
% and unesco2021recommendation.

\chapter{Discussion}

\section{Contributions of the internship}

This internship delivers a documented corpus of 1,153 French-language
resources on territorial adaptation to ageing, containing approximately
917,554 word tokens. It brings together heterogeneous material from six
institutional and associative sources: CNAV, Fondation de France, SIAGE,
HCFEA, IReSP Autonomie, and the French Network of Age-friendly Cities
(RFVAA). Beyond the corpus itself, the project provides a traceable data
architecture that preserves raw documents, intermediate files, structured
records, and review artefacts.

The work also proposes a practical natural language processing workflow for a
documentary context in which sources differ substantially in form and level of
detail. A transparent rule-based baseline is complemented by a local LLM-based
extractor with a constrained French prompt and JSON output. The resulting
records can be explored through fields such as target audience, initiative
type, theme, territorial scale, territory, holder, and date. A deterministic
normalisation layer then maps the extracted themes and initiative types to
controlled French vocabularies while retaining the original model outputs.

Finally, the internship makes the resulting index more auditable. Each
structured record is associated, where available, with provenance fields such
as the source site, source URL, original title, and collection date. This
distinction between extracted metadata and documentary provenance is important:
the former supports exploration, whereas the latter allows a user to return to
the source and verify an interpretation. The project should therefore be seen
as both an NLP prototype and a documented foundation for a future INSIGHT
index.

\section{Interpretation of the extraction results}

The human review of a stratified sample of 60 documents provides evidence that
the constrained extraction workflow is usable for several fields, but it does
not establish perfect reliability for the whole corpus. In particular, the
four controlled categorical fields obtained very high agreement with the gold
annotations: target audience, initiative type, and theme reached a micro-F1 of
1.000, while territorial scale reached a micro-F1 of 0.993. These results are
consistent with the fact that a restricted label set reduces output variation
and makes both automatic parsing and manual assessment easier.

The open-text fields require a more nuanced reading. On evaluable cases,
accepted rates were 0.981 for territory, 0.974 for initiative holder, and
0.786 for date. These figures should not be interpreted as an intrinsic
property of the model alone. They reflect the interaction between the model,
the prompt, the extraction limit, and--above all--the evidence available in the
source documents. A holder or date may be absent, implicit, or ambiguous in a
document. For this reason, the annotation protocol distinguished an incorrect
answer from a correct absence: when the source did not provide a reliable value
and the system returned an empty field, this was considered an accepted
outcome. Conversely, non-evaluable cases were excluded from the denominator
because the source did not allow a confident human decision either.

The source-level analysis further shows why an aggregate score is not
sufficient. The controlled fields were stable across the six subsets, whereas
the date field varied more strongly according to source format. For example,
dates were less often accepted in the CNAV and RFVAA subsets than in sources
whose documents expose publication dates more consistently. Territorial scale
also missed one label in the Fondation de France subset. These results do not
demonstrate that one source is inherently more difficult than another: each
subset contains only ten documents. They nevertheless identify concrete
priorities for improvement, particularly a clearer temporal schema and
source-specific checks for fields that are often incomplete.

Overall, the experiment supports the use of the metadata as an exploratory and
reviewable index. It does not justify presenting every extracted value as a
verified fact, nor does it replace human judgement for public-facing
publication or policy analysis. The evaluation is best understood as a first
benchmark that can be extended with further annotations, a second annotator,
and agreement measures between annotators.

\section{Interpretation of descriptive outputs}

The normalised records make it possible to describe the documentary collection
and to compare subsets of documents. The most frequent thematic labels include
\textit{autonomie et avanc\'ee en \^age}, \textit{sant\'e et pr\'evention}, and
\textit{lien social et lutte contre l'isolement}. Among initiative types,
workshops, training, studies or surveys, awareness-raising activities, and
events are recurrent. These patterns are plausible in a corpus devoted to
adaptation to ageing, but they must be read as properties of the collected and
labelled documents rather than as estimates of the real distribution of
initiatives in France.

Several mechanisms can influence these descriptive outputs. First, the corpus
is highly unbalanced across sources: RFVAA documents represent the largest
share of the collection. Second, the same initiative may be documented by more
than one resource, whereas other initiatives may have little or no web
presence. Third, source editors decide which activities to publish and how to
describe them. Finally, a single document can receive several themes or
initiative types; category counts are therefore not mutually exclusive.

The normalisation step improves comparability without eliminating these
limitations. It applies deterministic mappings to the original extracted
labels, so it does not create a second model prediction. However, any
controlled vocabulary necessarily groups together some distinctions present in
the source material. The preservation of the original fields, the use of an
\textit{autre / \`a revoir} category, and versioning of the taxonomy are thus
important safeguards. In this context, the figures and tables generated by the
project are most useful for exploratory navigation, hypothesis generation, and
discussion with INSIGHT users. They should not be used to rank territories,
assess local performance, or infer causal effects.

\section{Ethical, legal, and methodological considerations}

The corpus is composed of publicly accessible institutional and associative
resources, supplemented by SIAGE action sheets supplied within the internship
context. Its purpose is documentary analysis; the system does not profile
individuals and must not be used to make consequential decisions about people,
organisations, or territories. Public accessibility is not equivalent to an
unlimited right of reuse. Collection and dissemination must remain compatible
with the terms of access of the source websites, intellectual-property rights,
database rights where applicable, and data-protection obligations
\cite{cnil2024reuse}.

Provenance is consequently both a technical and an ethical requirement. When a
document comes from an external website, the index should retain its source
site and URL, as well as its original title. This makes attribution possible,
allows corrections to be traced back to their origin, and enables a user to
consult the source rather than treating the extracted metadata as a substitute
for it. In a future public interface, the index should preferentially expose
short metadata and links to the original resource, rather than redistributing
full documents without a clear legal basis.

The use of an LLM also requires calibrated claims. Structured outputs can be
wrong, incomplete, or overly confident even when they are syntactically valid.
The project mitigates this risk through controlled vocabularies, raw-output
archiving, provenance fields, and manual evaluation. These measures improve
auditability, but human review remains necessary before metadata are published
as factual statements. This is particularly important for territories, holders,
and dates, for which the source evidence can be ambiguous.

The project also has an environmental dimension. Local inference avoids
transmitting document content to a commercial remote API, but it does not make
the workflow impact-free. Energy use depends on the hardware, model, number of
requests, and repeated experimentation. As no energy measurement was performed,
the internship does not claim a quantified carbon footprint. Instead, it adopts
proportionate measures: a relatively small local model, no model training,
bounded input and output lengths, deterministic inference settings, and the
archiving of outputs to avoid unnecessary repeated runs. These choices are in
line with calls to consider the environmental costs of NLP and the broader
responsibility of AI systems \cite{strubell2019energy, unesco2021recommendation}.

\section{Limitations and future work}

The first limitation concerns corpus coverage. The collection is not a
representative sample of all initiatives related to ageing in France: it
inherits the selection and editorial practices of the six sources, and its
source distribution is strongly unbalanced. Some documents are web pages,
others are PDFs or action sheets; their textual quality, metadata, and layout
vary substantially. Web pages can also change after collection. The corpus
should therefore be described as a documented collection, not as an exhaustive
inventory of French initiatives.

The evaluation is deliberately modest. It uses 60 documents, ten per source,
and one human reviewer. This design is useful for detecting systematic
problems and comparing source formats, but it does not provide a large enough
sample for precise generalisation. Moreover, only the fields selected in the
annotation workbook were evaluated; free-text descriptions and every
normalised mapping were not independently assessed. Ten corpus documents do
not yet have an LLM output and are consequently absent from the structured
descriptive analyses. The next version should complete these records, enlarge
the benchmark, introduce double annotation on an overlapping subset, report
inter-annotator agreement, and assess the normalised fields directly.

Date extraction deserves specific attention. A single field can refer to a
publication date, an event date, a project start date, a duration, or no date
at all. A more explicit temporal schema, accompanied by source-specific rules
or evidence spans, would make the field more informative and easier to
evaluate. Similar improvements could be made for holder and territory by
storing the textual evidence that supports an extracted value. Future data
collection should also populate \texttt{collection\_date} at the time of
acquisition, rather than attempting to reconstruct it afterwards.

Further technical work could compare several local models and prompt versions,
add duplicate or near-duplicate detection, and develop automatic checks for
missing or contradictory metadata. Finally, a faceted search interface or a
hybrid lexical--semantic retrieval system could make the index more useful to
INSIGHT users. Such an interface should be designed with users, expose the
source and uncertainty of metadata, and keep human validation in the loop for
uses that go beyond exploratory documentary search.
